\begin{table*}[p]
\scriptsize
\centering
\caption{\textbf{Incremental contribution of scaffolding features in integrated adjacent-turn models.} Likelihood-ratio block tests compare nested logistic regressions for whether the next user turn is constructive. All models include prior user state, user framing, task ecology, assistant-turn index and the indicated dataset or model/source fixed effects.}\label{tab:integrated_scaffolding_block_tests}
\setlength{\tabcolsep}{4pt}
\renewcommand{\arraystretch}{1.08}
\resizebox{\textwidth}{!}{%
\begin{tabular}{llcc}
\toprule
Scope & Added feature block & LR $\chi^2$ (df) & p value \\
\midrule
six task settings, dataset FE & Broad scaffolded support after user/context controls & 584.6 (1) & $<.001$ \\
six task settings, dataset FE & M1--M6 descriptors within scaffolded support & 896.7 (6) & $<.001$ \\
six task settings, dataset FE & Full scaffolding block after user/context controls & 1481.3 (7) & $<.001$ \\
six task settings, model/source FE & Broad scaffolded support after user/context controls & 509.7 (1) & $<.001$ \\
six task settings, model/source FE & M1--M6 descriptors within scaffolded support & 829.8 (6) & $<.001$ \\
six task settings, model/source FE & Full scaffolding block after user/context controls & 1339.6 (7) & $<.001$ \\
WildChat only, model FE & Broad scaffolded support after user/context controls & 418.6 (1) & $<.001$ \\
WildChat only, model FE & M1--M6 descriptors within scaffolded support & 519.4 (6) & $<.001$ \\
WildChat only, model FE & Full scaffolding block after user/context controls & 937.9 (7) & $<.001$ \\
\bottomrule
\end{tabular}%
}
\end{table*}
